\documentclass{report}

\usepackage[ngerman]{babel}
\usepackage[utf8]{inputenc}
\usepackage[T1]{fontenc}
\usepackage{hyperref}
\usepackage{csquotes}
\usepackage[a4paper]{geometry}
\usepackage{graphicx}

\usepackage[
    backend=biber,
    style=apa,
    sortlocale=de_DE,
    natbib=true,
    url=false,
    doi=false,
    sortcites=true,
    sorting=nyt,
    isbn=false,
    hyperref=true,
    backref=false,
    giveninits=false,
    eprint=false]{biblatex}
\addbibresource{../references/bibliography.bib}


\title{Ethik im Umgang mit Daten}
\author{Dishan Selvanantham}
\date{\today}


\begin{document}

\maketitle

\abstract{
    In dieser Arbeit werden die Themen KI, Ethik und Daten behandelt, wie sie miteinander zusammenhängen, sowie welche Gefahren hinter der KI stecken.
}

\tableofcontents

\chapter{Einleitung}

\section{Was ist die KI}

Die KI, Abkürzung für Künstliche Intelligenz, ist ein Versuch menschliches Lernen und Intelligenz auf einen Computer zu übertragen und somit ihm Intelligenz zu verleihen. Diese Technologie ist in der Arbeitswelt bereits sehr weit verbreitet, da sie Prozesse optimieren und viele Aufgaben übernehmen kann und somit auch zeit spart. Es wird zwischen starker und schwacher KI unterschieden. Eine starke KI ist eine Maschine, die Probleme genereller Art lösen kann. Eine schwache KI hingegen bezeichnet eine Maschine, die nur spezielle Aufgaben bearbeiten kann und die Lösungswege vorher selbstständig erlernt. Auch wenn die KI im ersten Blick als nützliche Maschine erscheint, verbergen sich zahlreiche Gefahren hinter ihr.

\section{Daten}

Daten sind die Nahrung der Künstlichen Intelligenz. Die Künstliche Intelligenz wird mit Daten gefüttert, damit die Maschine lernen kann und Entscheidungen treffen kann. Die Qualität der Daten ist entscheidend für die Leistungsfähigkeit der KI. Unternehmen und Entwickler müssen aber sicherstellen, dass ihre KI- Anwendungen die Privatsphäre schützen und den ethischen Grundsätzen entsprechen. Der Umgang mit solch einer grossen Datenmenge muss also genau geregelt sein.

\section{KI-Ethik}

KI- Ethik ist ein Teilgebiet, das sich mit ethischen Fragen in Zusammenhang mit der Künstlichen Intelligenz befasst. Die Ethik in der KI soll sicherstellen, dass die Künstliche Intelligenz nur zum Wohle der Gesellschaft verwendet wird und keine Schäden verusacht. Die KI hat das Potenzial unser Leben zu revolutionieren und hat somit sehr viel Macht. Bei solch einer grossen Macht ist eine ebenfalls grosse Verantwortung notwendig.

\chapter{Funktionen}
\label{chap:methode}

\section{Wie funktioniert die KI}
Das Verhalten der Technologien werden generell mit Mitteln der Informatik und Mathematik simuliert. Die Maschinen werden für bestimmte Aufgaben trainiert, indem sie mit grossen Datenmengen gefüttert werden und das Muster erkennen. Die KI verwendet künstliche neuronale Netze, um zu funktionieren. Dabei wird die Funktionsweise des Gehirns nachgeahmt. Der Unterschied, das Gehirn kann viel mehr leisten als alle Computer. "Computer stellen diese Informationsverarbeitung  des Gehirns durch künstliche neuronale Netze nach. Informationen werden als Input auf der einen Seite eingegeben, verarbeitet und das Ergebnis wird auf der anderen Seite als Output wieder ausgegeben. Solche Systeme setzen sich aus sogenannten Algorithmen zusammen"\citep{ai-genius-community}.

\section{Algorithmus}
Ein Algorithmus kann als Anleitung betrachtet werden, das einem Computer sagt, was er machen muss, um eine Aufgabe oder Problem zu lösen. Also muss der Computer wissen, was wir von ihm wollen, bevor er die Aufgabe erledigen kann. Dies wird von IT-Spezialisten programmiert, die dem Computer erklären, was sie von ihm erwünschen und Anweisungen geben, was das System machen soll.

\section{Deep Learning}
"Deep Learning" (deutsch: mehrschichtiges Lernen oder tiefgehendes Lernen)ist eine Kategorie des maschinellen Lernens.
Im Prozess "Deep Learning"   wird dem Computer beigebracht, wie man komplexe Muster aus grossen Datenmengen verstehen kann, indem Schichten von künstlichen neuronalen Netzwerken verwendet werden. Die Konfiguration der Neuronen im Netzwerk basiert auf das Training mit grossen Datenmengen.

\section{Machine Learning}
Ein weiterer wichtiger Begriff im Bereich der Künstlicher Intelligenz ist das "Machine Learning" (deutsch: Maschinelles Lernen), die Begabung des Computers aus Erfahrungen zu lernen und Erkenntnisse aus Daten zu gewinnen, ohne dafür extra programmiert zu sein.

\chapter{Anwendungen der KI}
\label{chap:Anwendungen}

\section{KI-Anwendungen}
Die KI ist nicht mehr wegzudenken, da sie bereits heute in vielen Bereichen unseres Alltags präsent ist. Wer nur an Chat Gpt denkt, liegt hier komplett falsch. Oft kommen wir in Berührung mit der KI, ohne es wirklich wahrzunehmen. Ob Schachspiel gegen einen Computer, Besuch eines Online Shops, Interaktion mit Sprachassistenten oder Verwendung eines Übersetzungsdienstes. In all diesen Situationen kommen wir in Kontakt mit der Künstlichen Intelligenz. Von einer Schwachen KI spricht man, wenn die KI auf den Bereich beschränkt ist, indem sie eingesetzt wird, also nur in einem spezifischen Bereich ziemlich gut ist. Dies ist heutzutage in den meisten Situationen der Fall, da es "noch" keine starke KI gibt, die das menschliche Gehirn übertrifft und auf verschiedenen Bereichen verwendet werden kann. 



\section{Kann die KI gefährlich werden?}
Die KI, der grosse technologische Fortschritt der heutigen Zeit, ihre  Entwicklung  schreitet rasch voran und die Einsatzmöglichkeiten nehmen zu. Neben den vielen Vorteilen, besteht auch eine Gefahr, dass die KI eines Tages schlauer wird als Menschen. Steht also wirklich das Verderben der Menschheit durch die KI vor? Kann die KI aber tatsächlich Arbeitskräfte ersetzen? Die KI, ein Segen oder doch ein Fluch? Die KI wird in der Zukunft den Arbeitsalltag vieler Menschen ändern. Studien zeigen, dass Berufe wie z.B Mathematiker oder Programmierer am meisten bedroht sind. Das muss jedoch nicht unbedingt heissen, dass viele Arbeitskräfte ihren Beruf nicht mehr ausführen können, sondern dass sie von lästigen Aufgaben entlastet werden und mehr Zeit haben für Tätigkeiten, bei denen man nicht auf Menschen verzichten kann. Die KI kann zum Beispiel viele Bereiche der Medizin revolutionieren. Die Themen Rassismus und Sexismus kommen im Bereich der KI leider nicht selten vor, da KI- Systeme wie z.B Chat Bots aus Daten lernen, welche häufig nicht neutral sind, wiederspiegelt das System dies und kann diskriminierende Entscheidungen treffen. In den letzten Wochen und Monaten haben Nachrichten über Deep Fakes durch AI zugenommen. Die Künstliche Intelligenz erlaubt es uns Fotos, Videos als auch Audiodateien zu manipulieren. Viele Stars und Politiker wurdern bereits Opfer von den sogenannten Deep Fakes. All diese Szenarien zeigen, dass wir sehr verantwortungsbewusst mit der KI umgehen müssen und ihre Einsatzmöglichkeiten wie z.B das Erstellen von Deep Fakes beschränken müssen. Jedoch will ich noch betonen, dass sich hinter der KI einige Gefahren verbergen, aber auch sehr viele positive Aspekte. In den sozialen Medien wird die Gefahr der KI übertrieben dargestellt, was eher einer unrealistischen Fiktion entspricht



\chapter{Umgang mit Daten}
\label{chap:Daten}
\begin{figure}[h]
    \centering
    \includegraphics[width=0.5\textwidth]{KI2.jpg} 
    \caption{Umgang mit Daten}
    \label{fig:ai}
\end{figure}

Daten machen einen Grossteil aller Informationen aus und sind somit von entscheidender Bedeutung in der heutigen Zeit. Diese Daten können durch die KI automatisiert ausgewertet werden. Der Umgang mit Daten erfordert grundlegende Werte wie Transparenz, Gerechtigkeit, Fairness, Verantwortung, Datenschutz und Privatsphäre. Um einen verantwortungsvollen Umgang mit Daten sicherzustellen, müssen diese Werte von Menschen und Unternehmen wertgeschätzt werden. Es ist wichtig, dass Unternehmen sich der Bedeutung von Datenschutz und Privatsphäre bewusst sind und nötige Massnahmen ergreifen, um die Daten ihrer Kunden zu schützen. Der Umgang mit Daten spielt eine zentrale Rolle in der Gesellschaft und es ist wichtig, diesen verantwortungsvoll zu gestalten.



\printbibliography

\end{document}
